\documentclass[final]{beamer}

% ====================
% Packages
% ====================

\usepackage[T1]{fontenc}
% LuaLaTeX handles UTF-8 natively - luainputenc not needed
\usepackage{lmodern}
\usepackage[size=custom, width=122, height=91, scale=1.2]{beamerposter}
\usetheme{gemini}
\usepackage{graphicx}
\usepackage{booktabs}
\usepackage{tikz}
\usepackage{pgfplots}
\pgfplotsset{compat=1.14}
\usepackage{anyfontsize}
\usepackage{ragged2e}
\usepackage{xcolor}
\usepackage{enumitem}
\usepackage{caption}   % needed for \captionof inside beamer

% ====================
% Custom Colors (matching HTML palette)
% ====================

\definecolor{accent}{RGB}{196,80,0}       % --accent: #C45000
\definecolor{accentdark}{RGB}{139,50,0}   % --accent-dark: #8B3200
\definecolor{textmid}{RGB}{74,58,44}      % --text-mid
\definecolor{textdim}{RGB}{138,122,108}   % --text-dim

\captionsetup{font=small, labelfont=bf, textfont=it, skip=4pt}

% ====================
% Background Image
% ====================

\setbeamertemplate{background}{%
  \begin{tikzpicture}
    \useasboundingbox (0,0) rectangle (\paperwidth,\paperheight);
    \node[anchor=south west, inner sep=0] at (0,0)
      {\includegraphics[width=\paperwidth,height=\paperheight]{background.jpg}};
    \fill[white,opacity=0.88] (0,0) rectangle (\paperwidth,\paperheight);
  \end{tikzpicture}%
}

% ====================
% Lengths
% ====================

\newlength{\sepwidth}
\newlength{\colwidth}
\newlength{\midcolwidth}

\setlength{\sepwidth}{0.020\paperwidth}
\setlength{\colwidth}{0.295\paperwidth}
\setlength{\midcolwidth}{0.330\paperwidth}

\newcommand{\separatorcolumn}{\begin{column}{\sepwidth}\end{column}}

% ====================
% Research-question item  (NO TikZ -- uses a plain colored box)
% ====================

\newcommand{\rqitem}[2]{%
  \vspace{4pt}\noindent
  \fcolorbox{accent}{accent}{%
    \color{white}\bfseries\small\,#1\,}%
  \hspace{6pt}%
  \begin{minipage}[t]{0.87\linewidth}
    \small #2
  \end{minipage}\par\vspace{6pt}%
}

% ====================
% Title / Header
% ====================

\title{Active Solar Irradiance Observer \textcolor{accent}{(ASIO)}}

\author{%
  Evan Williams $\cdot$ Sarah Pawlowski $\cdot$ Chris Entzel $\cdot$
  Charles Kankelborg $\cdot$ Camrynn Wells $\cdot$ Caleb Pascual $\cdot$
  Reed Lando $\cdot$ Jake Davis $\cdot$ Chase Bronson $\cdot$
  Tyler Holliday $\cdot$ Emily Schwartz $\cdot$ Aidan Rasmussen $\cdot$
  Connor Parrott $\cdot$ Larry Springer $\cdot$ Ian Garcia $\cdot$
  Jacob Weber $\cdot$ Finnegan Davidson $\cdot$ Sophia Daly}

\institute[MSU]{%
  Montana State University $-$ Department of Physics $-$
  Space Science \& Engineering Laboratory}

% ====================
% Footer
% ====================

\footercontent{%
  MSU Student Research Celebration --- Spring 2026 \hfill
  Space Science \& Engineering Laboratory \hfill
  Montana State University $\cdot$ Department of Physics}

% ====================
% Logos
% ====================

\logoleft{%
  
\includegraphics[height=7cm]{msulogo.png}}

\logoright{%
  
\includegraphics[height=7cm]{nasalogo.png}\hspace{8pt}%
  
\includegraphics[height=5.5cm]{muselogo.png}\hspace{8pt}%
  
\includegraphics[height=5.5cm]{asiologo.png}}

% ====================
% Body
% ====================

\begin{document}

\begin{frame}[t]
\begin{columns}[t]
\separatorcolumn

% ============================================================
%  LEFT COLUMN
% ============================================================
\begin{column}{\colwidth}

  % -- Introduction ------------------------------------------
  \begin{block}{Introduction}

    Solar flares are among the most energetic events in our solar system,
    releasing enormous quantities of energy across the electromagnetic
    spectrum in timescales ranging from milliseconds to hours. The underlying
    physical mechanisms --- particularly sub-second oscillatory phenomena
    known as \textbf{\textcolor{accentdark}{Quasi-Periodic Pulsations
    (QPPs)}} --- remain poorly understood because existing solar monitors
    lack the temporal resolution and multi-band simultaneity needed to
    resolve them.

    \vspace{6pt}
    QPPs have been detected in microwave, hard X-ray, and EUV emissions, but
    the temporal relationships \textit{between} these bands at sub-second
    cadence are largely unconstrained. The spatial dynamics of the soft X-ray
    source during a flare --- and the thermal evolution of the flare late
    phase on centi-second timescales --- are similarly unexplored
    observational frontiers.

    \vspace{6pt}
    The \textbf{\textcolor{accentdark}{Active Solar Irradiance Observer
    (ASIO)}} is designed specifically to fill these gaps. Scheduled for
    launch in 2027 aboard NASA's \textbf{\textcolor{accentdark}{MUSE
    (MUlti-slit Solar Explorer)}} spacecraft, ASIO delivers simultaneous
    high-cadence irradiance measurements in the SXR and HXR bands with
    integrated coarse source localization --- a capability not available in
    any existing solar instrument.

  \end{block}

  % -- Research Objectives -----------------------------------
  \begin{block}{Research Objectives}

    \rqitem{1}{Do sub-second \textbf{\textcolor{accentdark}{Quasi-Periodic
      Pulsations (QPPs)}} exist simultaneously across the SXR, HXR, and EUV
      energy bands? Are there measurable temporal offsets between bands that
      reveal the causal sequence of energy release?}

    \rqitem{2}{Are \textbf{\textcolor{accentdark}{soft X-ray sources}} during
      solar flares spatially stationary, or do they exhibit rapid motion
      across the solar disk? If dynamic, over what spatial extent and on what
      timescale?}

    \rqitem{3}{How does the \textbf{\textcolor{accentdark}{thermal structure}}
      of a solar flare evolve during its late phase on centi-second timescales
      --- and what temperature and emission measure trajectories emerge as the
      plasma cools?}

    \vspace{6pt}
    \begin{columns}[t]
      \begin{column}{0.49\linewidth}
        \centering
        \includegraphics[width=\linewidth, height=5.5cm,
          keepaspectratio=false, clip]{slide2_img1.jpg}
        \captionof{figure}{\textbf{Fig.~1.} EUV observation of a solar
          flare. Scale markers indicate the source spans 50--190~Mm across
          the solar disk.}
      \end{column}
      \begin{column}{0.49\linewidth}
        \centering
        \includegraphics[width=\linewidth, height=5.5cm,
          keepaspectratio]{slide2_img2.png}
        \captionof{figure}{\textbf{Fig.~2.} QPP classification by period
          and duration (Tan et al.\ 2010). ASIO targets the fast and slow
          VSP regimes at sub-second cadence.}
      \end{column}
    \end{columns}

  \end{block}

  % -- References --------------------------------------------
  \begin{block}{References}

    \footnotesize
    \begin{enumerate}[leftmargin=18pt,
        label={\textcolor{accent}{\textbf{\arabic*.}}}]
      \item Ning, H., Chen, Y., Wu, Z., Su, Y., Tian, H., Li, G., Du, G.,
        \& Song, H. (2018). Two-stage energy release process of a confined
        flare with double HXR peaks. \textit{The Astrophysical Journal},
        \textit{854}(2), 178.
      \item De~Pontieu, B., Cheung, M., \& MUSE team. (2022). The
        MUlti-slit Solar Explorer (MUSE). \textit{Bulletin of the AAS},
        \textit{54}(7).
      \item Tan, B., Zhang, Y., Tan, C., \& Liu, Y. (2010). Microwave
        quasi-periodic pulsations in multi-timescales associated with a
        solar flare/CME event. \textit{The Astrophysical Journal},
        \textit{723}(1), 25--39.
    \end{enumerate}

  \end{block}

\end{column}

\separatorcolumn

% ============================================================
%  MIDDLE COLUMN
% ============================================================
\begin{column}{\midcolwidth}

  % -- The Instrument ----------------------------------------
  \begin{block}{The Instrument}

    ASIO is a compact photometer engineered as a hosted payload on the
    \textbf{\textcolor{accentdark}{NASA MUSE spacecraft}}, which will
    occupy a Sun-synchronous low Earth orbit at $\sim$667~km. The
    instrument leverages MUSE's spacecraft bus for power, thermal control,
    and telemetry while providing a high-cadence X-ray measurement channel
    fully independent of MUSE's primary UV spectrograph.

    \vspace{6pt}
    The flight chassis is precision-machined from aluminum with multiple
    aperture channels, each coupling solar X-rays to dedicated solid-state
    photodiode detectors. A \textbf{\textcolor{accentdark}{shadow plate}}
    upstream of the SXR detector array provides coarse source localization
    without imaging optics --- differential shadow patterns across detector
    segments encode the angular position of the X-ray source on the solar
    disk. An adjustable mounting ring permits fine alignment of the
    instrument axis after spacecraft integration.

    \vspace{8pt}
    \begin{columns}[t]
      \begin{column}{0.535\linewidth}
        \centering
        \includegraphics[width=\linewidth, height=5.4cm,
          keepaspectratio=false, clip]{slide3_img3.jpg}
        \captionof{figure}{\textbf{Fig.~3.} ASIO flight instrument ---
          precision-machined aluminum housing with multiple X-ray aperture
          channels and solid-state photodiode detectors.}
      \end{column}
      \begin{column}{0.445\linewidth}
        \centering
        \includegraphics[width=\linewidth, height=5.4cm,
          keepaspectratio]{slide3_img4.png}
        \captionof{figure}{\textbf{Fig.~4.} MUSE spacecraft CAD rendering.
          ASIO flies as a hosted payload on the main bus, launching
          in~2027.}
      \end{column}
    \end{columns}

    \vspace{6pt}
    \begin{columns}[t]
      \begin{column}{0.49\linewidth}
        \centering
        \includegraphics[width=\linewidth, height=4.5cm,
          keepaspectratio]{slide3_img5.png}
        \captionof{figure}{\textbf{Fig.~5.} ASIO structural CAD model
          showing aperture geometry and detector mounting.}
      \end{column}
      \begin{column}{0.49\linewidth}
        \centering
        \includegraphics[width=\linewidth, height=4.5cm,
          keepaspectratio]{slide3_img6.png}
        \captionof{figure}{\textbf{Fig.~6.} Adjustable positioning ring
          enabling precise instrument alignment after spacecraft
          integration.}
      \end{column}
    \end{columns}

  \end{block}

  % -- Alignment Procedure -----------------------------------
  \begin{block}{Alignment Procedure}

    \begin{columns}[t]
      \begin{column}{0.62\linewidth}
        The shadow plate upstream of the SXR photodiode array encodes the
        angular position of a solar X-ray source as a differential shadow
        pattern. Comparing flux ratios between detector segments yields a
        coarse source centroid without imaging optics. Post-integration
        alignment is characterised with an automated edge-detection
        pipeline: exterior and interior aperture corners are located to
        sub-pixel accuracy, compared against the optical model, and the
        adjustable ring is locked once verified.
      \end{column}
      \begin{column}{0.35\linewidth}
        \centering
        \includegraphics[width=\linewidth, height=5.6cm,
          keepaspectratio]{slide5_img10.png}
        \captionof{figure}{\textbf{Fig.~7.} Edge-detection alignment
          output.}
      \end{column}
    \end{columns}

  \end{block}

\end{column}

\separatorcolumn

% ============================================================
%  RIGHT COLUMN
% ============================================================
\begin{column}{\colwidth}

  % -- Instrument Testing ------------------------------------
  \begin{block}{Instrument Testing}

    All ASIO assembly and integration was performed under
    \textbf{\textcolor{accentdark}{Class~100 cleanroom conditions}} at
    MSU's SSEL facility. Calibrated red laser sources illuminated the
    aperture assembly on an optical bench to map the spatial response of
    each photodiode channel and validate the shadow-plate positional
    sensing concept. Differential illumination patterns matched ray-trace
    model predictions, confirming readiness for flight integration.

    \vspace{8pt}
    \begin{columns}[t]
      \begin{column}{0.32\linewidth}
        \centering
        \includegraphics[width=\linewidth, height=4.5cm,
          keepaspectratio=false, clip]{slide4_img7.jpg}
        \captionof{figure}{\textbf{Fig.~8.} ASIO integration in the
          Class~100 cleanroom facility at MSU~SSEL.}
      \end{column}
      \begin{column}{0.32\linewidth}
        \centering
        \includegraphics[width=\linewidth, height=4.5cm,
          keepaspectratio=false, clip,
          angle=270]{slide4_img9.jpg}
        \captionof{figure}{\textbf{Fig.~9.} ASIO on the optical bench
          during laser aperture characterization.}
      \end{column}
      \begin{column}{0.32\linewidth}
        \centering
        \includegraphics[width=\linewidth, height=4.5cm,
          keepaspectratio=false, clip]{slide4_img8.jpg}
        \captionof{figure}{\textbf{Fig.~10.} Red laser illumination of the
          shadow-plate aperture validating positional sensing geometry.}
      \end{column}
    \end{columns}

  \end{block}

  % -- Data Analysis -----------------------------------------
  \begin{block}{Data Analysis}

    ASIO's science pipeline converts multi-band irradiance time series into
    \textbf{\textcolor{accentdark}{temperature}} and
    \textbf{\textcolor{accentdark}{emission measure}} using ratios of SXR
    and HXR fluxes under an isothermal plasma assumption. Quasi-Periodic
    Pulsations are identified via \textbf{\textcolor{accentdark}{continuous
    wavelet decomposition}}, which resolves the time-frequency structure of
    the irradiance signal and isolates statistically significant periodic
    components above global (95\%) and local (99\%) significance
    thresholds. Synthetic signal validation --- shown below --- confirms
    robust QPP recovery from realistic noise environments prior to
    first-light solar observations.

    \vspace{8pt}
    \begin{columns}[t]
      \begin{column}{0.49\linewidth}
        \centering
        \includegraphics[width=\linewidth, height=3.3cm,
          keepaspectratio]{slide6_img11.png}
        \captionof{figure}{\textbf{Fig.~11.} Simulated irradiance time
          series with embedded quasi-periodic structure and additive noise
          --- the pipeline validation input.}
      \end{column}
      \begin{column}{0.49\linewidth}
        \centering
        \includegraphics[width=\linewidth, height=3.3cm,
          keepaspectratio=false, clip]{slide6_img12.png}
        \captionof{figure}{\textbf{Fig.~12.} Rectified wavelet power
          spectrum. QPP periods are clearly recovered above the 95\% global
          significance contour (red).}
      \end{column}
    \end{columns}

    \vspace{6pt}
    \centering
    \includegraphics[width=0.96\linewidth, height=4.7cm,
      keepaspectratio]{slide6_img13.png}
    \captionof{figure}{\textbf{Fig.~13.} Reconstruction from significant
      wavelet coefficients only (RMSE $= 1.86\times10^{-2}$, 6.7\% of
      std), demonstrating high-fidelity QPP isolation from background
      noise.}

  \end{block}

  % -- Acknowledgments ---------------------------------------
  \begin{block}{Acknowledgments}

    This work is supported by NASA under the MUSE Mission development
    program. Additional support provided by the Montana State University
    Space Science \& Engineering Laboratory (SSEL) and the MSU Department
    of Physics. Presented at the \textbf{MSU Student Research Celebration
    --- Spring 2026}.

  \end{block}

\end{column}

\separatorcolumn
\end{columns}
\end{frame}

\end{document}